\documentclass[10pt,a4paper]{article}
\usepackage[T1]{fontenc}
\usepackage{lmodern,microtype,amsmath,amssymb,booktabs,tabularx,enumitem}
\usepackage[margin=23mm]{geometry}
\usepackage[dvipsnames]{xcolor}
\usepackage[colorlinks=true,linkcolor=MidnightBlue,urlcolor=MidnightBlue,citecolor=MidnightBlue]{hyperref}
\usepackage{listings}
\lstset{basicstyle=\ttfamily\small,columns=fullflexible,keepspaces=true,
  frame=single,rulecolor=\color{black!20},breaklines=true}
\setlist{nosep,leftmargin=*}
\newcommand{\any}{\mathsf{Atom}}
\newcommand{\List}{\mathsf{ListOf}}
\newcommand{\checks}[3]{#1\vdash #2\Leftarrow #3}
\newcommand{\code}[1]{{\frenchspacing\texttt{\detokenize{#1}}}}
\newcommand{\name}[1]{\nolinkurl{#1}}
\hypersetup{pdftitle={Runtime Membership Contracts for SWI-PeTTa},
  pdfsubject={Membership, open constraints, implementation and proof boundary},
  pdfauthor={PeTTa runtime typing draft}}
\title{Runtime Membership Contracts for SWI-PeTTa}
\author{Research draft}
\date{6 November 2026}
\begin{document}
\maketitle
\begin{abstract}
PeTTa compiles MeTTa to SWI-Prolog. In the runtime typing draft described here,
a declaration such as \code{(: f (-> A B C))} is a membership contract: a call
of \code{f} is admitted only when its arguments belong to \code{A} and \code{B},
and its result must belong to \code{C}. Checking never runs an equation body a
second time and never binds a program variable. An argument that is still
unbound becomes a constraint that is checked when it binds. This note gives the
judgments, the reading of \code{Atom} and of type parameters, the generated
code, what the Lean model proves, and measurements against the current stable
branch: every typed workload measured costs fewer inferences and less CPU time.
It is a research draft, not a verified compiler or a static type system.
\end{abstract}

\section{Contracts on values}
PeTTa programs manipulate symbolic terms, unify variables, enumerate alternative
answers and edit spaces while they run. A type check at a function boundary
must therefore accept partially known data, preserve sharing, and fail rather
than add or repeat answers. Reading each declaration as a predicate on values
meets these requirements:
\begin{lstlisting}
(: Nil (List $a))
(: Cons (-> $a (List $a) (List $a)))
(: same (-> $a $a Bool))
(= (same $x $y) True)
\end{lstlisting}
A call \code{(same 1 2)} is admitted because both arguments are numbers, and
\code{(same 1 "s")} is refused because no single type contains both. The
repeated \code{$a} asks for one common type; it says nothing about how
\code{same} is implemented. Constructor parameters are fresh at each
constructor node, whereas the parameters of a function signature have one
assignment per call.

The design follows polymorphic logic programming~\cite{mycroft} and the
separation of goals from constraints in constraint logic programming~\cite{clp}.
Declarations guard execution at run time; they do not prove equation bodies
well typed in advance. SWI-Prolog supplies the mechanisms: indexed clauses for
membership tests, unification for private type equations, attributed variables
for requirements that wait for a value~\cite{attvar}, and \code{freeze/2} for
independent primitive tests~\cite{freeze}. Enforcement reads the declarations
directly. Reflection through \code{get-type} stays separate, so applications
may extend it without changing what a contract admits.

\section{Membership and type parameters}
Values $v$ are finite first-order terms; open values may also contain program
variables. $\Gamma$ is a declaration view. A type $p$ is built from
\code{Number}, \code{String}, \code{Bool}, \code{Symbol}, the structural forms
\code{ListOf}, \code{TupleOf} and \code{PairOf}, declared families such as
\code{(List $a)}, function types, and \code{Atom}, which inside a type is the
type of every value. A well-formed type is \emph{proper} if it does not
mention \code{Atom}. Write $\checks{\Gamma}{v}{p}$ when the closed value $v$ belongs to
the closed type $p$. Representative rules are:
\[
 \frac{v\text{ is represented as a number}}
      {\checks{\Gamma}{v}{\mathsf{Number}}}
 \qquad
 \frac{}{\checks{\Gamma}{v}{\any}}
 \qquad
 \frac{\checks{\Gamma}{v_i}{p}\quad(1\leq i\leq n)}
      {\checks{\Gamma}{[v_1,\ldots,v_n]}{\List(p)}}
\]
\[
 \frac{c:p_1,\ldots,p_n\to K(\bar q)\in\Gamma
       \quad \checks{\Gamma}{v_i}{\rho p_i}\ (1\leq i\leq n)}
      {\checks{\Gamma}{c(v_1,\ldots,v_n)}{K(\rho\bar q)}} .
\]
A nominal fact \code{(: c T)} gives membership of the whole value \code{c}; it
says nothing about fields. In the constructor rule, $\rho$ instantiates the
constructor's own parameters jointly, so repeated parameters stay equal.
Formation is a premise: \code{(ListOf ListOf)} is malformed and admits no value,
not even \code{()}. Primitive types test the representation; a declaration
cannot make a symbol a \code{Number}. An atom with no proper declaration is a
\code{Symbol}. Recursive families whose type grows at each step are checked by
the general relation instead of generated clauses.

\paragraph{Where Atom appears.}
\code{Atom} has different roles in a signature and inside a type:
\begin{center}
\begin{tabularx}{\linewidth}{@{}lX@{}}
\toprule
Position & Meaning\\ \midrule
Whole input slot, \code{(-> Atom ...)} & Pass the argument unevaluated; no value restriction.\\
Whole result slot, in every declaration & Return the equation body as a template.\\
Inside a structural type or field & Every value, so \code{(ListOf Atom)} holds \code{(1 "a" foo)}.\\
Index of a declared family & The index \code{Atom} itself; families are invariant.\\
\code{_}, \texttt{\%Undefined\%} & A whole slot with no restriction; inside a type they admit nothing.\\
\bottomrule
\end{tabularx}
\end{center}
Families are invariant because nothing declares how \code{(List $a)} varies with
\code{$a}: \code{(List Atom)} holds \code{(Cons 1 (Cons "a" Nil))}, built at
index \code{Atom}, but not a value declared \code{(List Number)}. Function types
are compared as written; there is no subtyping relation.

\paragraph{Parameters are proper.}
A signature parameter ranges over proper types only. Otherwise \code{Atom},
which every value inhabits, would solve every shared parameter, and
\code{(same 1 "s")} would pass. The rule holds inside compound types too: given
\code{(: wide (Index Atom Atom))} and \code{(: f (-> (Index $a $a) $a $a Bool))},
the call \code{(f wide 1 "s")} must not succeed by setting $\$a=\any$, and it
is refused in every argument order. The runtime attaches a delayed
proper-type obligation to each signature parameter. Two consequences follow. A
fact \code{(: x Atom)} is no evidence, so \code{x} is just a \code{Symbol}. And
a value declared \code{(List Atom)} meets a requirement \code{(List Atom)} but
not \code{(List $a)}.

\section{One call, one body computation}
For known inputs $\bar v$, a signature $d=(p_1,\ldots,p_n\to p_r)$ admits the
call with result $r$ when
\[
 \operatorname{Joint}_{d}(\Gamma_0,\Gamma_1,\bar v,r)
 \iff
 \exists\theta:\mathsf{Proper}.\
 \left(\bigwedge_i\checks{\Gamma_0}{v_i}{\theta p_i}\right)
 \land \checks{\Gamma_1}{r}{\theta p_r},
\]
where $\Gamma_0$ is read at entry and $\Gamma_1$ at return. Several declarations
for one function form a disjunction of such formulas. They are alternative
justifications for the same computation, so the body runs once. Entry keeps
every declaration the inputs admit, and the result may be justified by any of
them, not only by the first one found. The equations alone determine the order
and number of answers.

Each checked function compiles to an entry predicate around a private worker
that holds its equations. For \code{(: inc (-> Number Number))} the generated
entry tests are native Prolog, and an unbound argument gets a frozen test:
\begin{lstlisting}
inc(X, R) :-
    ( var(X) -> type_freeze(X, number(X)) ; number(X) ),
    '$body$9'(X, R),
    ( var(R) -> type_freeze(R, number(R)) ; number(R) ).
\end{lstlisting}
A shared parameter needs a frame that relates the arguments and the result.
For \code{same} the entry is
\code{types_enter([[T,T,'_']], [X,Y], F)}, then the worker, then
\code{types_exit(F, R)}, then the \code{Bool} test on \code{R}.
Argument modes are chosen before arguments are evaluated, and the arguments
captured by a partial application are not evaluated again. Direct, indirect,
partial and memoized calls all reach the same entry. The result is checked after
the body, so a check never precedes the body's effects or exceptions.

\paragraph{Open values.}
An unbound argument produces a residual requirement, never an invented value.
Type variables private to the check may be unified; program variables are not.
One residual formula keeps correlated arguments and alternative declarations
together. A binding fails when the evidence then known contradicts the formula;
otherwise the requirement waits for later bindings. The check does not decide
whether an open requirement can still be met.

\paragraph{Entry evidence and declaration edits.}
The parts of an input known at entry keep the evidence they had at entry; a part
that is still open records its evidence when it binds. Nothing else of the
declaration view is saved. If \code{old} and \code{new} both have type
\code{A}, a call \code{(keep old)} under \code{(: keep (-> $a $a))} may remove
\code{(: old A)} and return \code{new}: the input's entry evidence still gives
\code{A}, and the result is checked against the declarations at return. When a
function returns one of its inputs unchanged, the input's evidence is reused
only if SWI-Prolog's generation counters show that no declaration changed
during the call~\cite{transactions}.

\section{Implementation}
\begin{center}
\begin{tabularx}{\linewidth}{@{}lrX@{}}
\toprule
File & Lines & Responsibility\\ \midrule
\code{types.pl} & 383 & Declaration index, evidence, joint frames, proper
parameters, local structural checks.\\
\code{type_compilation.pl} & 234 & Entry and worker generation, argument modes,
tail and suffix reuse, retirement after edits.\\
\code{type_membership.pl} & 77 & Indexed clauses for closed membership, with
the general relation as fallback.\\
\code{type_obligations.pl} & 69 & Attributes on open values: evidence cells,
watchers, copying and release.\\
\bottomrule
\end{tabularx}
\end{center}
Line counts include comments. The old call-site checks are removed from the
translator, which shrinks by 46 lines. The first three files share the runtime
module, so generated names and the definition lifecycle are unchanged; the
attribute code is its own module.

A frame is \code{choices(Options, Stage, Identity, Stamp, Values)}, and an
empty option list means the frame is discharged. Each open program variable
holds one evidence cell and a list of watching frames without duplicates. An
input whose parameter occurs elsewhere only in the result constrains nothing
until the result exists; it only has to record later bindings. A call that is
still running already records every binding of the variables it watches, so a
nested call on the same open value does not watch them again, and recursion
on one open value takes linear rather than quadratic time. At return, anything
still unchecked arms its own watchers. All of this
state is undone by ordinary backtracking.

Answers are collected with ordinary \code{findall}: constraints on returned
variables are copied with them, and constraints on purely local variables are
dropped. Dropping them is sound only if those local constraints are
satisfiable. Adding an atom to a space stores plain variables, as SWI-Prolog's
assert does.

\section{What is proved}
The Lean library \name{Mettapedia.Logic.ConstraintTypes}~\cite{lean} states the
abstract model and proves, among others:
\begin{enumerate}
\item \textbf{Checker exactness.}
\name{DeclarationTyping.check_iff} and \name{OpenTyping.compile_correct}: the
executable checkers decide their declaration judgments.
\item \textbf{Staged exactness.}
\name{CallTyping.frame_exact} and \name{PatternContracts.frame_exact}: one
assignment serves both the input and the result conditions.
\item \textbf{Compiled membership.}
\name{DeterminedMembership.compiled_membership_correct}, under its premise that
fields determine membership; a nominal fact about a whole value is not
constructor evidence.
\item \textbf{Unshared parameters.}
\name{UnsharedParameters.waive_joint_iff}: a parameter that is an input's
whole type and occurs nowhere else may go unchecked when every value has some
proper type, which the canonical ground model provides. This is coverage, not
parametricity.
\item \textbf{Cost laws.}
\name{CheckCost.incremental_exact} and \name{entry_check_covers_suffixes}
justify per-part structural checks and suffix reuse;
\name{CheckingPatterns.matches_any} covers \code{Atom} inside a type.
\end{enumerate}
These are theorems about abstract finite-tree models. No proof yet connects
them to the source normalizer, the Prolog emitter, SWI's attributed variables
or the interpreter as a whole, so those are tested, not proved. Checks that
traverse structure reject cyclic terms; whole slots without a restriction do
not inspect them.

\section{Cost and compatibility}
Checking every suffix of an $n$-element list from scratch costs
$n(n+1)/2$ element checks. For structural types over primitives and
\code{Atom}, a suspension checks each newly bound cell once. A self-call on
\code{cdr-atom} of a list already checked against \code{(ListOf T)} skips that
slot; the other checks remain. Residual formulas for
nominal or overloaded types can still rescan values.

The table compares the same typed programs on both trees through compiled MeTTa
calls, including stable's call-site checks, as medians of fifteen interleaved
runs against \code{stable-fixed 5d4b75e}. The two complete examples include
loading and assertions. Inference counts are less sensitive to machine load
than CPU time.
\begin{center}
\begin{tabular}{@{}lrrrr@{}}
\toprule
& \multicolumn{2}{c}{Inferences} & \multicolumn{2}{c}{CPU seconds}\\
\cmidrule(lr){2-3}\cmidrule(l){4-5}
Workload & stable & draft & stable & draft\\ \midrule
10,000 scalar calls & 90,005 & 60,005 & .0026 & .0023\\
Counting, 200,000 steps & 2,000,014 & 1,000,009 & .088 & .016\\
Constructor walk, 400 cells & 1,851,823 & 83,012 & .159 & .0041\\
Open-variable loop & 90,014 & 70,025 & .0036 & .0016\\
\code{nilbc} & 20,100,081 & 19,486,151 & .583 & .533\\
\code{he_minimalmetta}, no declarations of its own & 24,199,749 & 20,844,814 & .755 & .662\\
\bottomrule
\end{tabular}
\end{center}
\code{nilbc} keeps its declaration \code{(: bc (-> $a Nat $b $b))}, whose
shared \code{$b} relates an open query to the answer; nested calls on one open
query are linear. The draft was faster in each of the 15 paired \code{nilbc}
runs. Ten small untyped programs have identical inference counts on both
trees, and three complete untyped examples stay within 0.5\%.

Behaviour changes relative to \code{stable-fixed} are deliberate:
\begin{center}
\begin{tabularx}{\linewidth}{@{}Xll@{}}
\toprule
Case & stable & draft\\ \midrule
\code{(: id (-> Number Number))} called through a variable on \code{"bad"} & \code{("bad")} & \code{()}\\
The same input bound to \code{"bad"} after the call & \code{("bad")} & \code{()}\\
One input with two nominal types & \code{(7 7)} & \code{(7)}\\
Body \code{(+ $x 1)} on 2, its \code{Atom} result declared after the equation & \code{3} & \code{(+ 2 1)}\\
Declarations \code{Number->Atom} and \code{String->Number}, body \code{(+ 1 2)}, input \code{"s"} & \code{()} & \code{(3)}\\
\code{(ListOf Atom)} input on \code{(1 "a" foo)} & \code{()} & \code{(7)}\\
\code{(List $a)} input on a value declared \code{(List Atom)} & \code{(7)} & \code{()}\\
\bottomrule
\end{tabularx}
\end{center}
In the mixed case an argument is held, and a result quoted, only when every
declaration says \code{Atom} at that position. Two example files change:
\code{typed_partial_arguments} no longer repeats answers and effects, and
\code{types_dependent} uses explicit guards for classifications derived through
\code{get-type}.

Validation: 113 runtime, 49 definition-order and 10 parser tests pass, normally
and under \code{-O}, and 177 example files with 786 assertions pass. Compiled and
general membership agree on 1,650 type/value pairs, for values given complete
and for values bound later, before and after declaration edits. The code is on
the \code{typing/minimal-runtime-stable-20261105} branch of
\href{https://github.com/godelclaw/PeTTa}{godelclaw/PeTTa}.

\begin{thebibliography}{9}
\bibitem{mycroft}
A. Mycroft and R. A. O'Keefe.
\emph{A polymorphic type system for Prolog}.
Artificial Intelligence, 23(3):295--307, 1984.
\href{https://doi.org/10.1016/0004-3702(84)90017-1}{doi:10.1016/0004-3702(84)90017-1}.
\bibitem{clp}
J. Jaffar and J.-L. Lassez.
\emph{Constraint Logic Programming}. POPL, 1987, 111--119.
\href{https://doi.org/10.1145/41625.41635}{doi:10.1145/41625.41635}.
\bibitem{attvar}
SWI-Prolog reference manual.
\href{https://www.swi-prolog.org/pldoc/man?section=attvar}{Attributed variables}.
\bibitem{freeze}
SWI-Prolog reference manual.
\href{https://www.swi-prolog.org/pldoc/doc_for?object=freeze/2}{\code{freeze/2}}.
\bibitem{transactions}
SWI-Prolog reference manual.
\href{https://www.swi-prolog.org/pldoc/man?section=transaction-impact}{Impact of transactions}.
\bibitem{lean}
MeTTapedia. Lean modules \name{Mettapedia.Logic.ConstraintTypes} and
\name{Mettapedia.Languages.MeTTa.PeTTa.RuntimeTyping}, commit \code{63cddbdaa9}.
\href{https://github.com/godelclaw/MeTTapedia}{github.com/godelclaw/MeTTapedia}.
\end{thebibliography}
\end{document}
